\section{Обсуждение и ограничения}
\label{sec:discussion}

Сопоставление рис.~\ref{fig:configs} и~\ref{fig:hierarchy} показывает,
почему оценка одной GEMM не заменяет проверку пути исполнения.
В исследованном стенде различия между вариантами GEMM ниже расхождения
с CPU-эталоном без XLA. Однако причины общего расхождения не разделены
между преобразованиями компилятора и внутренними форматами операций.

Результаты распространения зависят от выбранной метрики и точки её
вычисления. Нормы локальных ошибок, знаковая доля $b$ и корреляция $c$
не восстанавливают ни их совместную ковариацию, ни направления
после переноса. Поэтому наблюдаемую связь архитектуры с $R$ нельзя
превращать в причинный вывод о роли одного слоя без дополнительной проверки.
Не оценивались также спектральные нормы якобианов и величина
остатка $r_{\mathrm{lin}}$ в~\eqref{eq:accum}.

Испытания используют фиксированную инициализацию и отдельные компоненты,
а не обученные полноглубинные модели. Побитовая воспроизводимость
не даёт дисперсии по входам и начальным параметрам; доверительные
интервалы для $R$, $b$, $c$ не получены.
Сопоставимость локального расхождения cuBLAS с машинным эпсилоном
\texttt{FP32} требует проверки устойчивости соответствующих отношений,
но сама по себе не задаёт универсальной границы измеримости.

Изолированный и совместный обратные проходы исследуют численные градиенты,
но не траекторию оптимизатора. В частности, не проверены предположения
о компенсации ошибок по шагам обучения и о влиянии знаковой асимметрии
на итоговое качество модели. Выбор реализации для обучения требует
измерений сходимости на целевой задаче.

Результаты времени относятся к показаниям конкретного симулятора
и вычислительного узла. Они не являются оценкой полного времени процесса,
энергопотребления или скорости целевого ускорителя.
Близость к CPU-эталону также не устанавливает арифметической близости
к устройству: необходимое сравнение с реальным оборудованием
в этих испытаниях не выполнялось.
